\section{Laws of composition (internal). Subgroup (18-Oct-1994)}

Let $(S, \varphi)$, with $S \neq \emptyset$ and $\varphi: S \times S \to S$.

\begin{example}\leavevmode
\begin{itemize}
    \item $A = \text{given set}$, $(\mathcal{P}(A), \cup)$ or $(\mathcal{P}(A), \cap)$.
    \item $(\{f \mid f: [a, b] \to \mathbb{R}\}, +)$.
\end{itemize}
\end{example}

\begin{definition}
If $S$ is a set $S \neq \emptyset$, a map $\varphi: S \times S \to S$ is called an \textbf{internal law of composition} on $S$.
$$ S \times S \to S $$
$$ (x, y) \mapsto \varphi(x, y) \text{ is called the composite of } x \text{ with } y $$
\end{definition}

\textbf{Notations:} $x * y, x + y$ (additive), $x \cdot y$ (multiplicative). \\
Other notations for $\varphi(x, y)$: $x \circ y$, $x \top y$, $x \bot y$, $x \square y$.

\begin{definition}[Properties]
A law of composition $\varphi$ is:
\begin{itemize}
    \item \textbf{associative} if $\varphi(x, \varphi(y, z)) = \varphi(\varphi(x, y), z)$, $(\forall) x, y, z \in S$.
    \item \textbf{commutative} if $(\forall) x, y \in S \implies \varphi(x, y) = \varphi(y, x)$.
\end{itemize}
\end{definition}

\begin{definition}
An element $e \in S$ is called a \textbf{neutral element} (or identity element) for $\varphi$ if:
$$ (\forall) x \in S \text{ we have } \varphi(e, x) = \varphi(x, e) = x $$
\end{definition}

\begin{proposition}
If a law of composition $\varphi$ admits a neutral element, then it is unique.
\end{proposition}
\begin{proof}
Assume there are two neutral elements, $e$ and $e'$.
\begin{itemize}
    \item Considering $e$ as the neutral element: $e' \varphi e = e \varphi e' = e'$.
    \item Considering $e'$ as the neutral element: $e \varphi e' = e' \varphi e = e$.
\end{itemize}
Therefore, $e = e'$.
\end{proof}

\begin{example}
$(2\mathbb{N}, \cdot)$ does not have a neutral element. (The neutral element for multiplication would be $1$, but $1 \notin 2\mathbb{N}$).
\end{example}

\begin{definition}
Let $(S, \varphi)$ be a law of composition with a neutral element $e$. An element $x \in S$ is \textbf{invertible} (symmetrizable) with respect to $\varphi \iff (\exists) x' \in S$ such that:
$$ \varphi(x, x') = \varphi(x', x) = e $$
\end{definition}

\begin{proposition}
1) If $\varphi$ is associative and has a neutral element, and $x$ is invertible $\implies$ its inverse is unique. \\
2) If $x, y \in S$ are invertible, then $\varphi(x, y)$ is invertible and $(\varphi(x, y))' = \varphi(y', x')$.
\end{proposition}

\begin{proof}
1) Let $x', x'' \in S$ such that:
\begin{align*}
    x x' &= x' x = e \\
    x x'' &= x'' x = e
\end{align*}
$$ x'' = x'' e = x'' (x x') = (x'' x) x' = e x' = x' $$
\end{proof}

% Crossed-out example in the original manuscript:
% Ex: (\{ f:[a,b] \to \mathbb{R} \}, \circ)
% f(x) = 2x - 3 and x = 3/2 \notin [a,b] \implies f = invertible...

\begin{definition}
Let $(S, \cdot)$ with $\cdot$ being associative, $x_1, \dots, x_n \in S$.
$$ x_1 \cdot x_2 \cdot \dots \cdot x_n = \begin{cases} x_1, & n=1 \\ (x_1 \dots x_{n-1})x_n, & n \ge 2 \end{cases} = \prod_{i=1}^n x_i $$
\end{definition}

\begin{lemma}
Let $(S, \cdot)$ with $\cdot$ being associative. Then $(\forall) x_1, \dots, x_n \in S$, $n \in \mathbb{N}^*$, $m \in \mathbb{N}^*$, $1 \le m < n$:
$$ (x_1 \dots x_m)(x_{m+1} \dots x_n) = x_1 \dots x_n $$
\end{lemma}

\begin{proof}
By induction on $n$. \\
For $n=3$, the possible partitions are: \\
$m=1 \implies x_1(x_2 x_3)$ \\
$m=2 \implies (x_1 x_2)x_3$
\end{proof}

\begin{corollary}
If $(S, \cdot)$ is associative, $(\forall) x_1, \dots, x_n \in S$, $n \ge 3$ and $(\forall) m_1, \dots, m_r \in \mathbb{N}^*$ with $m_1 + \dots + m_r = n$, we have:
$$ (x_1 \dots x_{m_1})(x_{m_1+1} \dots x_{m_1+m_2}) \dots (x_{n-m_r+1} \dots x_n) = x_1 \dots x_n $$
\end{corollary}

\begin{example}
(Generalization of commutativity). Assume $(S, \cdot)$ is commutative. It can be shown that $(\forall) x_1, \dots, x_n \in S$, $n \ge 2$ and $(\forall) \sigma \in S_n$:
$$ x_{\sigma(1)} x_{\sigma(2)} \dots x_{\sigma(n)} = x_1 x_2 \dots x_n $$
\end{example}

\begin{example}\leavevmode
\begin{itemize}
    \item $(\mathbb{Z}, +), (\mathbb{Q}, +), (\mathbb{R}, +), (\mathbb{C}, +)$.
    \item $(\mathbb{Z}, \cdot), (\mathbb{Q}, \cdot), (\mathbb{R}, \cdot), (\mathbb{C}, \cdot)$.
    \item $(\mathcal{P}(A), \cup), (\mathcal{P}(A), \cap), (\mathcal{P}(A), \Delta)$ where $A \Delta B = (A \setminus B) \cup (B \setminus A)$. (The neutral element is $\emptyset$).
\end{itemize}
\end{example}

% Crossed-out example in the original manuscript:
% (\mathcal{P}(A), \setminus) \rightarrow A \setminus B = A \cap \complement_A B.
% associative ? A \setminus \emptyset = A, \emptyset \setminus A = \emptyset (It is not commutative).

\begin{definition}
Let $S \neq \emptyset$. A pair $(S, \varphi)$ where $\varphi$ is a law of composition on $S$ is called a \textbf{groupoid}.
\end{definition}

\begin{example}\leavevmode
\begin{itemize}
    \item $(\mathbb{N}, +), (\mathbb{N}, \cdot), (\mathbb{Z}, +)$ etc. are groupoids.
    \item $(\mathcal{M}_n(\mathbb{C}), +), (\mathcal{M}_n(\mathbb{C}), \cdot)$ are groupoids.
\end{itemize}
\end{example}

\begin{definition}
A groupoid $(S, \varphi)$ is called a \textbf{semigroup} if $\varphi$ is associative.
\end{definition}

\begin{example}
All of the above.
\end{example}

\begin{definition}
A semigroup which admits a neutral element is called a \textbf{monoid}.
\end{definition}

\begin{example}\leavevmode
\begin{itemize}
    \item $(\mathbb{N}, +) \rightarrow e = 0$; $(\mathbb{N}, \cdot) \rightarrow e = 1$ (similarly for $\mathbb{Z}, \mathbb{Q}, \mathbb{R}, \mathbb{C}$).
    \item $(\mathcal{M}_n(\mathbb{C}), +) \rightarrow e = O_n$; $(\mathcal{M}_n(\mathbb{C}), \cdot) \rightarrow e = I_n$.
    \item $(\mathcal{P}(A), \cup) \rightarrow e = \emptyset$; $(\mathcal{P}(A), \cap) \rightarrow e = A$; $(\mathcal{P}(A), \Delta) \rightarrow e = \emptyset$.
    \item $(\mathcal{P}(A), \setminus)$ - no (it does not have a neutral element).
    \item $(A^A, \circ)$ where $A^A = \{f \mid f: A \rightarrow A\}$, the neutral element is $1_A$.
\end{itemize}
\end{example}

\vspace{1em}
\textbf{Calculation rules in a semigroup}

Let $(S, \cdot)$ be a semigroup and $x \in S$.
\begin{align*}
    x^1 &= x \\
    x^2 &= x \cdot x \\
    &\dots \\
    x^n &= \underbrace{x \cdot x \dots x}_{n \text{ times}}
\end{align*}
$$ x^n \cdot x^m = x^{n+m} $$
$$ (x^n)^m = x^{n \cdot m} $$
If $S$ is a monoid: $x^0 \stackrel{\text{def}}{=} e$ (by convention).

\textit{Additive form:}
\begin{align*}
    1 \cdot x &= x \\
    2x &= x + x \\
    &\dots \\
    nx &= \underbrace{x + x \dots x}_{n \text{ times}}
\end{align*}
$$ nx + mx = (n+m)x $$
$$ m(nx) = (mn)x $$
If $S$ is a monoid: $0 \cdot x = 0$.